% HD-F1-T25. GPT-6 Astra Pro; /root/transport; 22 November 2026 UTC.
\section[HD-F1-T25: the residual Gaussian source has a rational double base]{HD-F1-T25: exactly two odd scalar inertias, an actual rational double-base coefficient, and the bounds $n\ge22$, $m\ge209$}
\label{sec:f1-gaussian-specht-double-base}

Retain the original natural sixfold source in the remaining
Gaussian alternating branch of T24. Its actual coefficient is
\[
 V=\mathbb Q(i)\otimes_{\mathbb Q}W_\lambda,
 \qquad \rho(g)=c(g)R(\sigma(g)),\qquad m=\dim W_\lambda,
 \tag{F1T25.1}
\]
where $\lambda\ne\lambda'$, $R$ is the rational Specht
representation, and the actual image is one of the three
subgroups of $\mu_4\times S_n$ proved in T24. In particular
$c(G)=\mu_4$. Write $g_1,\ldots,g_b$ for the original
product-one generating inertia tuple and $a_j=\operatorname{rank}
(\rho(g_j)-1)$. The carried source equalities are
\[
 b\ge12,\qquad \sum_j a_j=2m+6,
 \qquad \operatorname{rank}(\rho(g)-1)\ge5
                  \quad(g\ne1).
 \tag{F1T25.2}
\]
All ranks below refer to this faithful original representation.
We prove that precisely two $c(g_j)$ equal $i$ or $-i$.
The actual quotient by $\ker(c^2)$ is a projective line.
Its rational coefficient and cup pairing are explicitly
identified below. Moreover
\[
                         n\ge22,\qquad m\ge209.
 \tag{F1T25.3}
\]
No Hermitian determinant class is inferred merely from the
double quotient; the inherited primitive corner and all pairing
scalars are retained.

\paragraph{1. Exactly two odd phases, including their actual squares.}
The eigenvalue multiplicities of the rational matrix
$R(\sigma(g))$ at $i$ and $-i$ agree. For an odd-phase
element $g$, let this multiplicity be $u$. Then
\[
 \operatorname{rank}(\rho(g)-1)=m-u\ge m/2,
 \qquad
 \operatorname{rank}(\rho(g)^2-1)=m-2u
              =2\operatorname{rank}(\rho(g)-1)-m.
 \tag{F1T25.4}
\]
Indeed $\rho(g)^2=-R(\sigma(g)^2)$, whose fixed space
is exactly the sum of the two $\pm i$ eigenspaces of
$R(\sigma(g))$.

The number $q$ of odd-phase inertia entries is even and
positive: apply $c^2$ to the product-one tuple and use
generation of $c(G)$. For $q\ge4$, the lower bound
$m\ge27$ and (F1T25.2) give
\[
 \sum_j a_j\ge q m/2+5(b-q)
       \ge 2m+5(b-4)\ge2m+40,
\]
a contradiction. Thus $q=2$.

Each of these two squares is a nonidentity original
group element. To verify this, $R$ is faithful on $S_n$:
its kernel intersects the simple $A_n$ trivially because
the restriction is nontrivial irreducible, and the remaining
normal-kernel possibility is also trivial. If
$R(\sigma(g)^2)=-1$, that permutation acts scalarly,
so faithfulness makes it central in $S_n$, hence the identity,
which acts by $+1$. This contradiction proves the claim.
Consequently (F1T25.2) applies to $g^2$, and
\[
  a_j\ge\left\lceil\frac{m+5}{2}\right\rceil
       \quad\text{at the two odd phases},\qquad
 2m+6\ge2\left\lceil\frac{m+5}{2}\right\rceil+50.
 \tag{F1T25.5}
\]
In particular $m\ge49$, with $m\ge50$ when $m$ is even.
The use of rank at least five for the squares is justified
by the actual nonidentity check above.

\paragraph{2. The quotient curve, every lifted inertia, and the rational coefficient.}
Let $C\to B=\mathbb P^1$ be the original full $G$-cover,
put $H=\ker(c^2)$, and form the actual quotient
\[
                   B'=C/H\longrightarrow C/G=B.
 \tag{F1T25.6}
\]
It is connected of degree two and ramifies precisely at
the two odd-phase branch values. Riemann--Hurwitz gives
$2g(B')-2=2(-2)+2=-2$, so $B'\simeq\mathbb P^1$.
For a nonodd original inertia, there are two points above
its branch value, with the two conjugate $H$-inertia
operators. For an odd inertia $g_j$, there is one point
above it with $H$-inertia $g_j^2$. These are the local
maps of the normalized quotient, equivalently the usual
local coordinates $z\mapsto z^e$ followed by its degree-two
intermediate quotient. The two odd squares are nontrivial
by paragraph 1, and all the other lifted operators are
nontrivial by faithfulness. Therefore the actual $H$-cover
over $B'$ has precisely $2b-2\ge22$ coefficient branch
points and total coefficient rank
\[
 2\sum_{c(g_j)^2=1}a_j+
       \sum_{c(g_j)^2=-1}(2a_j-m)=2m+12.
 \tag{F1T25.7}
\]

Define the rational $H$-module $U=W_\lambda$ by
$h\cdot u=c(h)R(\sigma(h))u$; here $c(h)=\pm1$.
Its restriction to $A_n$ is the given absolutely
irreducible rational module, so $U$ is absolutely
irreducible. Choose an actual $t\in G$ with $c(t)=i$.
The exact induction map is
\[
 \begin{split}
 \Phi_{\mathbb Q}:\mathbb Q[G]\otimes_{\mathbb Q[H]}U
       &\longrightarrow\operatorname{Res}_{K/\mathbb Q}V,\\
 g\otimes u&\longmapsto\rho(g)u.
 \end{split}
 \tag{F1T25.8}
\]
The $H$-balancing relation proves well-definedness, and
left multiplication proves $G$-equivariance. On the two
cosets $1,t$, it is $(v,w)\mapsto v+iR(\sigma(t))w$.
Its inverse is obtained by taking rational real and
imaginary parts and applying $R(\sigma(t))^{-1}$ to
the latter. After extension to $K$, the equivalent map is
\[
 (v,w)\longmapsto
     (v+iR(\sigma(t))w,\ v-iR(\sigma(t))w)
          \in V\oplus\overline V.
 \tag{F1T25.9}
\]
Complex conjugation exchanges the two components, so this
is exactly the scalar extension of (F1T25.8).

Choose the invariant positive rational symmetric form $b$
on $W_\lambda$ by rational averaging. The original invariant
Hermitian coefficient form is a positive rational multiple
of its extension, by absolute irreducibility; absorb that
specified multiple in $b$. On the cosets $1,t$ the induced
rational coefficient form is $b\perp b$, and
\[
 \Phi_{\mathbb Q}^*\operatorname{Tr}_{K/\mathbb Q}h
                         =2(b\perp b).
 \tag{F1T25.10}
\]
This follows by expanding $v+iR(t)w$ and using the
orthogonality of $R(t)$. Multiplication by $i$ transported
to the induced module is the explicit rational operator
\[
                I(v,w)=(-R(t)w,R(t)^{-1}v),
          \qquad I^2=-1,\qquad I^\dagger=-I.
 \tag{F1T25.11}
\]

Over the original branch complement, induction is the
pushforward of the coefficient system $U$ on the actual
double cover. On a lifted cellular decomposition, the
cochain map collects the two sheet values into the two
coset summands. It commutes with the cellular differentials,
with restriction to boundary circles, and with the
absolute-to-relative cohomology map. At an odd branch
circle it identifies the invariant space with the
invariants of $g_j^2$; at a nonodd circle it is the sum
of the two conjugate invariant spaces. Thus it extends
to the actual parabolic images. Cap product is unchanged
under the sheet-collecting map because the induced pairing
is the sum over those two sheets. Applying (F1T25.8) then
multiplies the rational pairing by exactly two, as in
(F1T25.10). This proves the parabolic identification and
its polarization scalar, rather than only a dimension
identity. Its rational parabolic dimension is twelve,
as also follows from (F1T25.7).

\paragraph{3. Literal primitive projectors and the original corner.}
For a nonzero rational coefficient vector $w_0\in U$,
put $q_0=b(w_0,w_0)$. The rational $H$ matrix-coefficient
projector is
\[
 e_0=\frac{m}{|H|q_0}
       \sum_{h\in H}b(w_0,\rho(h^{-1})w_0)h.
 \tag{F1T25.12}
\]
Schur orthogonality makes it the primitive self-adjoint
projector onto the coefficient line $\mathbb Qw_0$.
The paired Gaussian projector in $\mathbb QG$ is
\[
 \frac{2m}{|G|q_0}\sum_{g\in G}
       \operatorname{Re}h(w_0,\rho(g^{-1})w_0)g.
 \tag{F1T25.13}
\]
For $g\notin H$ its coefficient is zero, since the
matrix coefficient is purely imaginary; for $g\in H$
it is exactly (F1T25.12), since $|G|=2|H|$.
Thus (F1T25.12) and (F1T25.13) are literally the same
element of $\mathbb QG$. The induction isomorphism
(F1T25.8) also shows that no other rational $G$ block
contains $U$. Hence this is a primitive projector in
the original Gaussian block, with corner $K$.
After clearing its denominator, its image is an actual
abelian subvariety of the original $JC$ with its
restricted polarization.

The prescribed original self-adjoint primitive projector
need not have a rational coefficient vector. Let its
$K$-line be $Kw$, with $q=h(w,w)>0$. In the original
matrix algebra component, the actual rank-one operator
\[
             A(x)=w_0\,h(w,x)/q
                   \in e_0\mathbb QG e
 \tag{F1T25.14}
\]
satisfies $A(w)=w_0$ and $A^\dagger A=(q_0/q)e$.
The membership follows from the already established
full matrix component over $K$; all its entries lie
in $K$. Choose an integer $d>0$ clearing its group-ring
coefficients. The actual homomorphism induced by $dA$
between the two connected images in $JC$ is an isogeny:
its rational homology map is invertible on the paired
coefficient line. Pullback/norm adjunction gives
\[
              (dA)^*\theta_{e_0}=d^2(q_0/q)\theta_e.
 \tag{F1T25.15}
\]
It intertwines the central Gaussian action. This is the
explicit comparison with the prescribed corner; the
positive scalar is not discarded. In rank six its
determinant contribution is the norm of
$(d^2q_0/q)^3$, but no value for the remaining determinant
is concluded.

\paragraph{4. An exact partition inequality from the original tuple.}
For the boxes $(r,s)$ of $\lambda$, set
\[
 p_1=\sum_{(r,s)\in\lambda}(s-r),\qquad
 p_2=\sum_{(r,s)\in\lambda}(s-r)^2.
\]
The three-cycle content formula and the square of the
transposition class sum give respectively
\[
 \begin{split}
 r_3:=\frac{\chi_\lambda(3)}m
   &=\frac{3p_2-3\binom n2}{n(n-1)(n-2)},\\
 r_{22}:=\frac{\chi_\lambda(2,2)}m
   &=\frac{4(p_1^2-3p_2+n(n-1))}
                 {n(n-1)(n-2)(n-3)}.
 \end{split}
 \tag{F1T25.16}
\]
For the second formula, the class-sum square is
$\binom n2\,1+3C_3+2C_{22}$, while its scalar on
$W_\lambda$ is $p_1^2$. These are the exact content
and class-sum identities read and proved at their
applied scope in T22; self-conjugacy is not used in
(F1T25.16).

For any nonidentity permutation $s$, choose a
transposition $\tau$ that does not commute with it.
Then $[s,\tau]$ is a three-cycle or double transposition.
The original rational Specht matrices give the exact
operator comparison, for every scalar $z$,
\[
 \operatorname{rank}(R([s,\tau])-1)
            \le2\operatorname{rank}(zR(s)-1).
\]
Rationality pairs the two nontrivial three-cycle
eigenvalues, so their moving rank is
$2(m-\chi_\lambda(3))/3$. A double transposition has
moving rank $(m-\chi_\lambda(2,2))/2$. Consequently
every nonidentity actual coefficient element has rank
at least
\[
 d_\lambda=
 \max\left\{5,
  \left\lceil\min\left(
       \frac{m(1-r_3)}3,\frac{m(1-r_{22})}4
                         \right)\right\rceil\right\}.
 \tag{F1T25.17}
\]
For scalar elements the rank is $m$, also satisfying
this bound. Apply it both to the ten or more nonodd
inertias and to the two nonidentity odd squares.
The original Euler equality yields the explicit
necessary inequality
\[
 2\left\lceil\frac{m+d_\lambda}{2}\right\rceil
                  +10d_\lambda\le2m+6.
 \tag{F1T25.18}
\]

For $\lambda=(n-2,2)$ or its conjugate, strengthen
$d_\lambda$ to $\max(d_\lambda,n-3)$; for
$\lambda=(n-2,1,1)$ or its conjugate, strengthen it
to $\max(d_\lambda,n-2)$. Here is the precise input:
Guralnick--Tiep, JEMS 14 (2012), Lemma 3.3 and its
complete proof, pp.~625--626, show for these two
characters and every nonidentity $s$ that
$|\chi_\lambda(s)|\le\chi_\lambda((12))$ for
$n\ge9$. Conjugating the partition changes a character
only by the sign, so the absolute bound persists.
For a finite-order operator, an eigenspace of
codimension $a$ gives $|\chi(s)|\ge m-2a$.
The displayed character formulas give
$(m-\chi((12)))/2=n-3,n-2$, respectively. This proves
the asserted all-scalar rank bounds on the actual
operators. No low-age hypothesis from other parts of
that paper is used.

\paragraph{5. The finite range and the coefficient-rank bound.}
The companion exact certificate enumerates every
partition for $7\le n\le21$. It computes $m$ both
by the Young-tableau corner recurrence and by the
Frobenius dimension formula, with an additional hook
formula equality check. For every non-self-conjugate
partition with $m\ge49$, the left side of
(F1T25.18), using the just-proved strengthened ranks,
exceeds $2m+6$. The number of such partitions and the
least positive difference at each $n$ are
\[
\begin{array}{c|rrrrrrrrrrrrrr}
n&8&9&10&11&12&13&14&15&16&17&18&19&20&21\\ \hline
\#&6&16&30&46&70&94&128&168&222&288&376&480&616&780\\
\min&92&104&84&116&40&40&38&36&34&30&24&18&12&4
\end{array}
\tag{F1T25.19}
\]
There is no eligible partition at $n=7$. Every row,
partition, dimension, content sum, bound, and difference
is retained in the certificate table. The recursion
enumerates a partition by its first part and the
remaining partition with bounded largest part, so it
is exhaustive. Thus $n\ge22$.

The standard and sign-standard shapes are already
excluded by the actual HD-G26 morphisms. To check the
identification literally, the permutation module on
tabloids of shape $(n-1,1)$ is $\mathbb Q^n$ by the
label in its singleton row; its Young-symmetrizer
submodule is the span of all $e_a-e_b$, exactly the
augmentation space. The sign-standard isomorphism is
the actual sign-twist Young-symmetrizer map given in
T22. It replaces the genuine scalar character by
$c\,\operatorname{sgn}\sigma$ and preserves its
correlations. All three actual T24 images satisfy
the exact subgroup hypotheses of G26, which allows
dependent scalar and permutation projections.
Its primitive-projector and polarization comparisons
therefore apply to the prescribed source. A period
rank-nine source of either shape is already split,
so these shapes do not belong to the remaining gate.

For completeness, the elementary dimension bound needed
now can be obtained directly from the same tableau
recurrence. For $n\ge9$, every shape other than the
one-dimensional and standard shapes and their conjugates
has dimension at least $n(n-3)/2$. The cases $n=9,10,11$
are explicitly checked in the certificate. Inductively,
a nonrectangular such shape has at least two removable
corners, none giving a one-dimensional shape. At least
one child is nonstandard: otherwise adding a box to a
standard child gives only a standard shape or
$(n-2,2),(n-2,1,1)$ and their conjugates, and each of
the latter has a nonstandard child. One child therefore
has dimension at least $(n-1)(n-4)/2$ and another at
least $n-2$, whose sum is $n(n-3)/2$.
For a rectangular shape $(a^b)$ with $a,b\ge2$, remove
its unique corner and then the two corners of the
result. For $n\ge12$ the two resulting shapes of size
$n-2$ are both nonstandard and not one-dimensional.
Their dimensions sum to at least
$(n-2)(n-5)\ge n(n-3)/2$. This proves the induction.
Applying it at $n\ge22$ gives $m\ge209$.

\paragraph{Scope of the advance.}
The conclusion restricts the actual residual Gaussian
alternating source, including its correlated scalar image,
to $n\ge22$, $m\ge209$, and (F1T25.18), with precisely
two odd scalar inertias. Its rational double-base
coefficient, complex structure, parabolic pairing,
primitive projector, and comparison to the prescribed
corner are the displayed maps. The result creates no
additional detector family and assigns no new determinant
class. In the condition-A source programme these are
restrictions on the same finite-cover correspondence;
the condition-I cycle requirement and the Cayley--Hamilton
division map still act on the original target through
the recorded isogeny and its exact polarization scalar.
